% !TEX program = xelatex
\documentclass[11pt,a4paper]{article}

% ---------- Encoding & Fonts ----------
\usepackage{fontspec}
\defaultfontfeatures{Ligatures=TeX}
\setmainfont{TeX Gyre Pagella} % 近似优雅衬线；可改为 Times New Roman
\setsansfont{TeX Gyre Heros}
\setmonofont{JetBrains Mono}[
  Scale=0.9,
  Contextuals=Alternate
]

% ---------- Geometry & Spacing ----------
\usepackage[a4paper,left=1.9cm,right=1.9cm,top=1.4cm,bottom=1.3cm,footskip=0.5cm]{geometry}
\usepackage{setspace}
\setstretch{1.05}
\usepackage{parskip}
\setlength{\parskip}{2pt}
\setlength{\parindent}{0pt}

% ---------- Colors ----------
\usepackage{xcolor}
\definecolor{secgray}{HTML}{000000}   % 原来 222222
\definecolor{lightgray}{HTML}{000000} % 原来 777777
\definecolor{linkblue}{HTML}{1A73E8}


% ---------- Links ----------
\usepackage[colorlinks=true, urlcolor=linkblue, linkcolor=black, citecolor=black]{hyperref}

% ---------- Page Style ----------
\usepackage{fancyhdr}
\pagestyle{fancy}
\fancyhf{} % 清除所有页眉页脚
\fancyfoot[C]{\thepage} % 居中显示页码
\renewcommand{\headrulewidth}{0pt} % 移除页眉线
\renewcommand{\footrulewidth}{0pt} % 移除页脚线

% ---------- Icons ----------
\usepackage{fontawesome5}

% ---------- Lists ----------
\usepackage{enumitem}
\setlist[itemize]{leftmargin=1.4em, itemsep=2pt, topsep=2pt, label={\small$\bullet$}}
\setlist[enumerate]{leftmargin=1.6em, itemsep=2pt, topsep=2pt}

% ---------- Titles ----------
\usepackage{titlesec}
\titleformat{\section}
  {\large\bfseries\color{secgray}\uppercase}
  {}{0pt}{}
\titlespacing*{\section}{0pt}{10pt}{6pt}

% ---------- Tabular helpers ----------
\usepackage{tabularx}
\newcolumntype{L}{>{\raggedright\arraybackslash}X}

% ----- Skills table helpers -----
\newlength{\skilllabelwd}
\setlength{\skilllabelwd}{2cm} % ← 左列宽度，按你版面微调 2.4–3.0cm
\newlength{\paperlabelwd}
\setlength{\paperlabelwd}{1cm} % ← 论文编号列宽度

\newcommand{\skillrow}[2]{%
  \noindent\begingroup\setlength{\tabcolsep}{0pt}%
  \begin{tabularx}{\linewidth}{@{}p{\skilllabelwd} X@{}}%
    #1 & #2%
  \end{tabularx}\par\endgroup
}

\newcommand{\paperrow}[2]{%
  \vspace{4pt}%
  \noindent\begingroup\setlength{\tabcolsep}{0pt}%
  \begin{tabularx}{\linewidth}{@{}p{\paperlabelwd} X@{}}%
    #1 & #2%
  \end{tabularx}\par\endgroup
}

% ---------- CV Header ----------
\newcommand{\cvname}[1]{%
  \begin{center}\LARGE\bfseries #1\end{center}
}
\newcommand{\cvcontact}[1]{%
  \begin{center}#1\end{center}
}
\newcommand{\cvaddress}[1]{%
  \begin{center}\small #1\end{center}
}
\newcommand{\cvupdated}[1]{%
  \begin{center}\footnotesize Last updated: #1\end{center}
}

% ---------- Entry Macros ----------
% left=主标题/单位 | right=日期/标注 | body=条目列表或文字（自动顶格）
\newcommand{\cventry}[3]{%
  \vspace{4pt}
  \noindent\begin{tabularx}{\linewidth}{@{}L r@{}}
    #1 & \textit{#2}
  \end{tabularx}\par
  #3
  \vspace{4pt}
}

% Research experience with bold titles
\newcommand{\resentry}[3]{%
  \vspace{4pt}
  \noindent\begin{tabularx}{\linewidth}{@{}L r@{}}
    \textbf{#1} & \textit{#2}
  \end{tabularx}\par
  #3
  \vspace{4pt}
}



% 教育 / 经验条目（带位置/导师）
% title=学位/职位, org=学校/机构, place=地点, date=日期, extra=附加行（导师/证书等，可留空）
\newcommand{\eduentry}[5]{%
  \vspace{4pt}
  \noindent\begin{tabularx}{\linewidth}{@{}L r@{}}
    \textbf{#1}, \textit{#2} & \textit{#4}
  \end{tabularx}\par
  \noindent{\small #3}\par
  \if\relax\detokenize{#5}\relax\else{\small #5\par}\fi
}

% Publications 编号项
% tag=[C1]/[J2]/[PP1] 等；正文可带\textit{}等
\newcommand{\pubitem}[2]{%
  \paperrow{\textbf{#1}}{#2}
}

% 小分组标题（Awards 下的 Paper Awards / Fellowships 等）
\newcommand{\subsec}[1]{\vspace{4pt}\textbf{#1}\par\vspace{2pt}}

% ---------- Document ----------
\begin{document}

% ===== Header =====
\cvname{Qiuyi Ding}
\cvcontact{%
  \faEnvelope\ \href{mailto:dingqy@umich.edu}{dingqy@umich.edu}\ \quad
  \faLink\ \href{http://qiuyiding.com}{http://qiuyiding.com}\ \quad
  \faGoogle\ \href{https://scholar.google.com/citations?hl=en&user=sn4ACXUAAAAJ}{Qiuyi Ding}
}
\cvaddress{University of Michigan, Ann Arbor, MI, USA}
\cvupdated{Nov 2025}

% ===== Research Interest =====
\section{Research Interest}
My research interests lie in controllable generation, multimodal perception, and unsupervised learning.

% ===== Education =====
\section{Education}
\eduentry{B.S., Computer Science}{University of Michigan}{}{2023 -- Present}{}

% ===== Research Experience =====
\section{Research Experience}
\cventry{\textbf{HidingSound: Natural Soundscape Generation for Online Noise Coverage}}{2025}{%
\textit{Advised by Prof. Andrew Owens}\par
\begin{itemize}
\item An online soundscape generation system using DiT to produce immersive natural soundscapes that perceptually mask unwanted noise in real-world environments
\end{itemize}
}
\cventry{\textbf{Representation-Guided Optimization for Diffusion and Vision-Language Models}}{2025}{%
\begin{itemize}
\item Improve the resolution and recognition capability of diffusion models and VLMs by optimizing in representation space, guided by DINO features as a perceptual signal for finer-grained semantic alignment
\end{itemize}
}
\cventry{\textbf{EXP-Bench: Can AI Conduct AI Research Experiments?}}{2025}{%
\textit{Advisor: Prof. Mosharaf Chowdhury, Prof. Ang Chen}\par
\begin{itemize}
\item We built Exp-Bench, a scalable benchmark for LLM evaluation by automating domain document curation, reasoning-oriented QA generation, and multi-agent verification, yielding diverse, high-fidelity datasets that better capture real-world complexity and align with human judgments
\end{itemize}
}
\cventry{\textbf{Curie: Automated and Rigorous Scientific Experimentation with AI Agents}}{2024}{%
\textit{Advisor: Prof. Mosharaf Chowdhury, Prof. Ang Chen}\par
\begin{itemize}
\item Curie is the first AI-agent framework designed for automated and rigorous scientific experimentation. Curie helps answer your curiosity through end-to-end experimentation automation, ensuring that every step—from hypothesis formulation to result interpretation—is conducted with precision, reliability, and reproducibility
\item \href{https://www.just-curieous.com}{Website}
\end{itemize}
}

% ===== Publications =====
\section{Publications \ [\texorpdfstring{\href{https://scholar.google.com/citations?hl=en&user=sn4ACXUAAAAJ}{\faGoogle}}{Google Scholar}]}
\textit{* equal contribution}

\pubitem{[1]}{\textcolor{linkblue}{\textbf{Hiding Sound Using Diffusion Illusions}}. \textbf{Qiuyi Ding}, Yiming Dou, Andrew Owens. \textit{(Honor Thesis, to be submitted)}.}
\pubitem{[2]}{\href{https://arxiv.org/abs/2505.24785}{\textbf{EXP-Bench: Can AI Conduct AI Research Experiments?}}. Patrick Tser Jern Kon, \textbf{Qiuyi Ding}, Jiachen Liu, Xinyi Zhu, Jingjia Peng, Jiarong Xing, Yibo Huang, Yiming Qiu, Jayanth Srinivasa, Myungjin Lee, Mosharaf Chowdhury, Matei Zaharia, Ang Chen. 2025. In: \textit{Proceedings of the International Conference on Learning Representations (ICLR 2026)}.}
\pubitem{[3]}{\href{https://arxiv.org/abs/2502.16069}{\textbf{Curie: Toward Rigorous and Automated Scientific Experimentation with AI Agents}}. Patrick Tser Jern Kon*, Jiachen Liu*, \textbf{Qiuyi Ding}, Yiming Qiu, Zhenning Yang, Yibo Huang, Jayanth Srinivasa, Myungjin Lee, Mosharaf Chowdhury, Ang Chen. 2025. In: \textit{arXiv preprint (2025)}.}

% ===== Projects =====
\section{Selected Projects}
\cventry{\textbf{RAG Benchmark for Pathology Education Assistant}}{2025}{%
\textit{Multidisciplinary Design Program, DIAG Team of UMich}\par
\begin{itemize}
\item Built a RAG-based benchmark for pathology education, combining LLMs with open educational resources for accurate, domain-adaptive knowledge support. \href{https://peer.asee.org/board-92-wip-generative-ai-based-learning-tutor-for-biomedical-data-science-gail-tutor-bds}{ASEE Annual Conference Poster}
\end{itemize}
}
\cventry{\textbf{Immersive NPCs in VR Game Environment}}{2024}{%
\textit{Course Project for EECS498: AI-Enabled Mixed Reality}\par
\begin{itemize}
\item Built LLM-powered NPCs in Unity + VRChat with GPT-4, giving them persistent memory, distinct personalities, and full speech I/O. NPCs that actually remember you next time.
\end{itemize}
}
\cventry{\textbf{Autonomous Hazardous Object Detection and Navigation}}{2025}{%
\textit{Course Project for EECS465: Introduction to Robotics}\par
\begin{itemize}
\item Built a robotics pipeline for autonomous hazardous object detection and path planning, using a VLM for open-vocabulary object recognition, 3D segmentation for spatial localization, and SLAM for real-time mapping and obstacle-aware navigation
\end{itemize}
}
\cventry{\textbf{Encoder Integration and Efficiency Optimization for Language-Embedded Radiance Fields}}{2024}{%
\textit{Course Project for EECS442: Computer Vision}\par
\begin{itemize}
\item LERF embeds CLIP language features into a NeRF, enabling open-vocabulary text queries to localize objects in 3D scenes. Swapped in a SigLIP encoder for sharper localization accuracy, and cut training time by 31\% via AdamW and reduced eval/checkpoint overhead.
\end{itemize}
}

% ===== Services =====
\section{Services}

\cventry{EECS245 Tutor (Mathematics for Machine Learning)}{2025}{}

% ===== Selected Awards & Honors =====
\section{Selected Awards \& Honors}

\cventry{University of Michigan Honors Program}{2026}{}
\cventry{James B. Angell Scholars}{2025}{}

% ===== Skills =====
\section{Skills}

\subsec{Programming Languages}
\skillrow{Proficient}{Python, C/C++, CUDA, Java, \TeX, HTML/CSS}
\skillrow{Familiar}{JAX, C\#, MATLAB, Mathematica, JavaScript}
\skillrow{Capable}{Bash, SQL, Triton, Pascal}

\subsec{Natural Languages}
\skillrow{Native}{Mandarin Chinese}
\skillrow{Fluent}{English}



\end{document}

